\documentclass[a4paper,12pt]{article}

\usepackage[T2A]{fontenc}
\usepackage[utf8]{inputenc}
\usepackage[russian]{babel}
\usepackage{amsmath,amssymb}
\usepackage{paratype}
\renewcommand{\familydefault}{\sfdefault}
\usepackage{icomma}
\usepackage[a4paper,margin=20mm,headheight=25pt,headsep=7mm,footskip=12mm]{geometry}
\usepackage{microtype}
\usepackage{tikz}
\usetikzlibrary{calc,arrows.meta,positioning,shapes.geometric}
\usepackage{tabularx,booktabs}
\usepackage{enumitem}
\usepackage{hyperref}
\usepackage{parskip}
\usepackage{fancyhdr}
\usepackage{setspace}
\usepackage{xcolor}

\definecolor{accent}{HTML}{008080}
\definecolor{darkaccent}{HTML}{004D4D}
\definecolor{textgray}{HTML}{2F3437}
\definecolor{softgray}{HTML}{667277}
\definecolor{rulegray}{HTML}{D5DEDE}

\hypersetup{
  colorlinks=true,
  linkcolor=darkaccent,
  urlcolor=darkaccent,
  pdfauthor={Артём Александрович Лёвин},
  pdftitle={Тригонометрические уравнения — чек-лист для Ярослава Гаврилова}
}

\color{textgray}
\onehalfspacing
\setlength{\parindent}{0pt}
\setlength{\emergencystretch}{2em}
\setlist[itemize]{leftmargin=1.45em,itemsep=0.3em,topsep=0.35em}
\setlist[enumerate]{leftmargin=1.75em,itemsep=0.45em,topsep=0.35em}
\renewcommand{\arraystretch}{1.25}

\pagestyle{fancy}
\fancyhf{}
\fancyhead[L]{\small\color{softgray}Ярослав Гаврилов · ЕГЭ}
\fancyhead[R]{\small\color{softgray}Тригонометрические уравнения}
\fancyfoot[C]{\small\color{softgray}\thepage}
\renewcommand{\headrulewidth}{0.3pt}
\renewcommand{\headrule}{%
  \hbox to\headwidth{\color{rulegray}\leaders\hrule height \headrulewidth\hfill}%
}

\fancypagestyle{flowpage}{%
  \fancyhf{}%
  \fancyfoot[C]{\small\color{softgray}\thepage}%
  \renewcommand{\headrulewidth}{0pt}%
  \renewcommand{\footrulewidth}{0pt}%
}

\newcommand{\sect}[1]{%
  \vspace{0.45em}%
  {\Large\bfseries\color{darkaccent}#1}\par
  \vspace{0.2em}%
  \noindent{\color{rulegray}\rule{\linewidth}{0.6pt}}%
  \vspace{0.35em}%
}

\newcommand{\subsect}[1]{%
  \vspace{0.35em}%
  {\large\bfseries\color{accent}#1}\par
  \vspace{0.1em}%
}

\newcommand{\key}[1]{\textbf{\color{darkaccent}#1}}
\newcommand{\Z}{\mathbb{Z}}
\newcommand{\R}{\mathbb{R}}

\tikzset{
  flowstart/.style={
    ellipse,
    draw=darkaccent,
    line width=0.9pt,
    align=center,
    minimum width=49mm,
    minimum height=11mm,
    fill=white,
    font=\small
  },
  flowaction/.style={
    rounded corners=3pt,
    draw=accent,
    line width=0.9pt,
    align=center,
    text width=39mm,
    minimum height=13mm,
    fill=white,
    inner sep=4pt,
    font=\small
  },
  flowwide/.style={
    rounded corners=3pt,
    draw=accent,
    line width=0.9pt,
    align=center,
    text width=100mm,
    minimum height=13mm,
    fill=white,
    inner sep=4pt,
    font=\small
  },
  flowdecision/.style={
    diamond,
    aspect=2.15,
    draw=darkaccent,
    line width=0.9pt,
    align=center,
    text width=27mm,
    inner sep=2pt,
    fill=white,
    font=\small
  },
  flowarrow/.style={
    -{Stealth[length=3mm,width=2mm]},
    line width=0.9pt,
    draw=textgray
  },
  flowlabel/.style={
    font=\small,
    fill=white,
    inner sep=1.1pt,
    text=textgray
  }
}

\begin{document}

\hypersetup{pageanchor=false}

\begin{titlepage}
\thispagestyle{empty}
\centering

\vspace*{23mm}

{\Large\color{softgray}Чек-лист}\par

\vspace{7mm}

{\Huge\bfseries\color{darkaccent}
Тригонометрические уравнения}\par

\vspace{4mm}

{\LARGE\bfseries\color{accent}
Преобразования · однородные уравнения · отбор корней}\par

\vspace{19mm}



{\normalsize 24 сентября 2026 г.}\par

\vspace{5mm}

{\normalsize\bfseries
Лёвин Артём Александрович эксклюзивно для Ярослава Гаврилова}\par

\vspace{18mm}

\begin{tikzpicture}[remember picture,overlay]
  \draw[accent!55,line width=1.1pt]
    ($(current page.south west)+(16mm,18mm)$)
    .. controls
    ($(current page.south west)+(59mm,31mm)$)
    and
    ($(current page.south east)+(-64mm,8mm)$)
    ..
    ($(current page.south east)+(-16mm,20mm)$);
\end{tikzpicture}

\end{titlepage}

\hypersetup{pageanchor=true}
\pagenumbering{arabic}

\sect{1. Главная траектория решения}

\begin{enumerate}
  \item
  \key{Сначала приведите аргументы к удобному виду.}
  Если рядом встречаются $x$ и $2x$, проверьте формулы двойного угла:
  \[
    \sin 2x=2\sin x\cos x,
  \]
  \[
    \cos 2x
    =2\cos^2x-1
    =1-2\sin^2x
    =\cos^2x-\sin^2x.
  \]

  \item
  \key{После каждого преобразования ищите общий множитель.}
  Если он появился, вынесите его за скобки.

  \item
  \key{Если получилось четыре слагаемых, проверьте группировку.}
  Часто удобно объединить их попарно так, чтобы после первого вынесения
  появились одинаковые скобки. Затем выполняется второе вынесение.

  \item
  \key{Не делите на выражение, которое может быть равно нулю.}
  Сначала обоснуйте, что делитель не обращается в ноль для корней
  уравнения, либо отдельно разберите нулевой случай.

  \item
  Если получено уравнение вида
  \[
    a\sin x+b\cos x=0,
  \]
  стремитесь перейти к одной тригонометрической функции.

  \item
  Если после общего решения требуется выбрать корни на промежутке,
  переводите условие для $x$ в условие для целого параметра $k$.
\end{enumerate}

\subsect{Полезная эвристика}

Количество слагаемых помогает увидеть структуру, но не заменяет анализ
выражения. Четыре слагаемых действительно часто ведут к группировке;
три слагаемых нередко удаётся упростить тождеством или заменой.
Способ выбирают по форме конкретного уравнения.

\sect{2. Формулы, которые нужны в этой теме}

\subsect{2.1. Простейшие тригонометрические уравнения}

\[
\sin x=a,\quad |a|\le1
\quad\Longrightarrow\quad
x=(-1)^k\arcsin a+\pi k,
\qquad k\in\Z.
\]

\[
\cos x=a,\quad |a|\le1
\quad\Longrightarrow\quad
x=\pm\arccos a+2\pi k,
\qquad k\in\Z.
\]

\[
\tg x=a
\quad\Longrightarrow\quad
x=\arctg a+\pi k,
\qquad k\in\Z.
\]

Для табличных значений записывайте точный угол. Например,
\[
\tg x=-1
\quad\Longrightarrow\quad
x=-\frac{\pi}{4}+\pi k,
\qquad k\in\Z.
\]

\subsect{2.2. Почему нельзя бездумно делить на переменное выражение}

Из равенства
\[
\sin x\,(2\cos x+1)=0
\]
нельзя сразу делить на $\sin x$: при $\sin x=0$ существуют корни,
которые после такого деления исчезнут.

Правильно использовать свойство нулевого произведения:
\[
\sin x=0
\quad\text{или}\quad
2\cos x+1=0.
\]

\key{Общее правило.}
Деление на выражение $F(x)$ является равносильным преобразованием только при
$F(x)\ne0$. Если это заранее не известно, сначала исследуйте случай
$F(x)=0$ или докажите, что он невозможен для корней данного уравнения.

\sect{3. Вынесение общего множителя}

Рассмотрим уравнение
\[
\sin 2x+\sin x=0.
\]

Сначала выравниваем аргументы:
\[
2\sin x\cos x+\sin x=0.
\]

Теперь виден общий множитель $\sin x$:
\[
\sin x(2\cos x+1)=0.
\]

По свойству нулевого произведения:
\[
\sin x=0
\quad\text{или}\quad
\cos x=-\frac12.
\]

Следовательно,
\[
x=\pi k
\quad\text{или}\quad
x=\pm\frac{2\pi}{3}+2\pi k,
\qquad k\in\Z.
\]

\key{Контрольный вопрос после каждой новой строки:}
«Можно ли сейчас вынести общий множитель?»

\sect{4. Группировка и двойное вынесение}

Алгебраическая модель приёма:
\[
AB+AC-DB-DC
=A(B+C)-D(B+C)
=(B+C)(A-D).
\]

Тригонометрический пример:
\[
2\sin x\cos x+\sin x-2\cos x-1=0.
\]

Группируем первое слагаемое со вторым, третье --- с четвёртым:
\[
(2\sin x\cos x+\sin x)+(-2\cos x-1)=0.
\]

Выполняем первое вынесение:
\[
\sin x(2\cos x+1)-(2\cos x+1)=0.
\]

Одинаковая скобка стала общим множителем:
\[
(2\cos x+1)(\sin x-1)=0.
\]

Отсюда
\[
\cos x=-\frac12
\quad\text{или}\quad
\sin x=1,
\]
то есть
\[
x=\pm\frac{2\pi}{3}+2\pi k
\quad\text{или}\quad
x=\frac{\pi}{2}+2\pi k,
\qquad k\in\Z.
\]

\key{Откуда появляется единица.}
Например,
\[
\sin x(2\cos x+1)-(2\cos x+1)
=(2\cos x+1)(\sin x-1),
\]
потому что второй множитель фактически имеет коэффициент $-1$.

\sect{5. Однородные уравнения первой степени}

Уравнение
\[
a\sin x+b\cos x=0
\]
называют однородным тригонометрическим уравнением первой степени относительно $\sin x$ и $\cos x$.

В рамках занятия основной приём --- переход к тангенсу делением
на $\cos x$.

Если $a\ne0$, корень такого уравнения не может удовлетворять условию
$\cos x=0$: тогда $\sin x=\pm1$, и левая часть равна
$\pm a\ne0$.

Следовательно, для корней
\[
\cos x\ne0,
\]
поэтому деление допустимо:
\[
a\tg x+b=0.
\]

Получаем
\[
\tg x=-\frac{b}{a},
\]
а значит,
\[
x=\arctg\!\left(-\frac{b}{a}\right)+\pi k,
\qquad k\in\Z.
\]

Если $a=0$ и $b\ne0$, исходное уравнение имеет вид
\[
b\cos x=0
\]
и решается как простейшее.

Если $a=b=0$, равенство выполняется при любом
\[
x\in\R.
\]

\subsect{Пример}

\[
3\sin x-2\cos x=0.
\]

Для корней $\cos x\ne0$, поэтому
\[
3\tg x-2=0,
\qquad
\tg x=\frac23.
\]

Ответ:
\[
x=\arctg\frac23+\pi k,
\qquad k\in\Z.
\]

\sect{6. Отбор корней на промежутке}

Пусть общее решение имеет вид
\[
x=\frac{\pi}{6}+\pi k,
\qquad k\in\Z,
\]
и требуется выбрать корни на промежутке
\[
-\frac{\pi}{3}\le x\le\frac{\pi}{2}.
\]

Подставляем формулу корней вместо $x$:
\[
-\frac{\pi}{3}
\le
\frac{\pi}{6}+\pi k
\le
\frac{\pi}{2}.
\]

Вычитаем $\frac{\pi}{6}$ из всех трёх частей, затем делим на $\pi>0$:
\[
-\frac{\pi}{2}\le\pi k\le\frac{\pi}{3}
\quad\Longrightarrow\quad
-\frac12\le k\le\frac13.
\]

Поскольку $k\in\Z$, получаем $k=0$, поэтому
\[
x=\frac{\pi}{6}.
\]

\begin{center}
\begin{tikzpicture}[x=2.2cm,y=1cm]
  \draw[-{Stealth[length=2.6mm]},line width=0.85pt]
    (-2.15,0)--(2.15,0)
    node[right] {$k$};

  \foreach \x/\lab in {
    -2/$-2$,
    -1/$-1$,
    0/$0$,
    1/$1$,
    2/$2$
  }{
    \draw[line width=0.7pt]
      (\x,0.08)--(\x,-0.08)
      node[below=3pt] {\lab};
  }

  \draw[accent,line width=2.1pt]
    (-0.5,0.33)--(0.333,0.33);

  \fill[accent] (-0.5,0.33) circle (1.7pt);
  \fill[accent] (0.333,0.33) circle (1.7pt);
  \fill[darkaccent] (0,0) circle (2.3pt);

  \node[above=3pt,text=accent]
    at (-0.5,0.33) {$-\frac12$};

  \node[above=3pt,text=accent]
    at (0.333,0.33) {$\frac13$};
\end{tikzpicture}
\end{center}

\key{Смысл параметра $k$.}
Он нумерует значения общего решения, возникающие из периодичности
тригонометрических функций. Поэтому при отборе корней сначала находят
допустимые целые $k$, затем соответствующие им значения $x$.

\clearpage
\thispagestyle{flowpage}

\begin{center}
{\Large\bfseries\color{darkaccent}
Алгоритм: выбор основного приёма решения}\par

\vspace{5mm}

\begin{tikzpicture}[node distance=5mm and 7mm]

  \node[flowstart] (start)
    {Тригонометрическое уравнение};

  \node[flowwide,below=of start] (prep)
    {Приведите аргументы к удобному виду и перенесите все слагаемые
    в одну часть};

  \node[flowdecision,below=5mm of prep] (factorq)
    {Есть общий множитель?};

  \node[flowaction,right=8mm of factorq] (factor)
    {Вынесите его; не потеряйте нулевой случай множителя};

  \node[flowdecision,below=5mm of factorq] (groupq)
    {Есть удачная группировка?};

  \node[flowaction,right=8mm of groupq] (group)
    {Сгруппируйте и выполните двойное вынесение};

  \node[flowdecision,below=5mm of groupq] (homq)
    {Получилось $a\sin x+b\cos x=0$?};

  \node[flowaction,right=8mm of homq] (hom)
    {При $a\ne0$ обоснуйте $\cos x\ne0$ и перейдите к $\tg x$};

  \node[flowwide,below=6mm of homq] (retry)
    {Иной вид: продолжите равносильные преобразования
    и снова оцените структуру};

  \node[flowwide,below=6mm of retry] (simple)
    {Решите полученные простейшие уравнения
    и запишите общее решение};

  \node[flowdecision,below=5mm of simple] (intervalq)
    {Нужен отбор корней?};

  \node[flowwide,below=6mm of intervalq] (select)
    {Подставьте формулу корней в границы промежутка,
    найдите целые $k$, затем вычислите соответствующие $x$};

  \node[flowstart,below=6mm of select] (finish)
    {Проверьте ограничения и запишите ответ};

  \draw[flowarrow]
    (start)--(prep);

  \draw[flowarrow]
    (prep)--(factorq);

  \draw[flowarrow]
    (factorq.east)
    -- node[flowlabel,above] {да}
    (factor.west);

  \draw[flowarrow]
    (factorq.south)
    -- node[flowlabel,right] {нет}
    (groupq.north);

  \draw[flowarrow]
    (groupq.east)
    -- node[flowlabel,above] {да}
    (group.west);

  \draw[flowarrow]
    (groupq.south)
    -- node[flowlabel,right] {нет}
    (homq.north);

  \draw[flowarrow]
    (homq.east)
    -- node[flowlabel,above] {да}
    (hom.west);

  \draw[flowarrow]
    (homq.south)
    -- node[flowlabel,right] {нет}
    (retry.north);

  \draw[flowarrow]
    (retry.west)
    -- ++(-10mm,0)
    |- (prep.west);

  \draw[flowarrow]
    (factor.east)
    -- ++(8mm,0)
    |- (simple.east);

  \draw[flowarrow]
    (group.east)
    -- ++(8mm,0)
    |- (simple.east);

  \draw[flowarrow]
    (hom.east)
    -- ++(8mm,0)
    |- (simple.east);

  \draw[flowarrow]
    (simple)--(intervalq);

  \draw[flowarrow]
    (intervalq.south)
    -- node[flowlabel,right] {да}
    (select.north);

  \draw[flowarrow]
    (select)--(finish);

  \draw[flowarrow]
    (intervalq.east)
    -- ++(34mm,0)
    |- node[flowlabel,pos=0.18,right] {нет}
    (finish.east);

\end{tikzpicture}
\end{center}

\clearpage

\sect{7. Типичные ошибки и самопроверка}

\begin{itemize}
  \item
  Потеря корней после деления на $\sin x$, $\cos x$ или другое
  выражение без проверки, может ли оно быть равно нулю.

  \item
  Пропуск формулы двойного угла, когда в одном уравнении встречаются
  аргументы $x$ и $2x$.

  \item
  Ошибка знака при вынесении минуса:
  \[
  -(A+B)=-A-B,
  \qquad
  -(A-B)=-A+B.
  \]

  \item
  Неудачная группировка. Её цель --- получить одинаковый множитель
  в двух группах и выполнить второе вынесение.

  \item
  Запись периода $2\pi k$ в уравнении с тангенсом.
  Для
  \[
  \tg x=a
  \]
  используется период $\pi$, поэтому в ответе стоит $\pi k$.

  \item
  Перебор значений $x$ при отборе корней.
  Надёжнее сначала получить интервал для $k$ и воспользоваться
  условием $k\in\Z$.

  \item
  Ошибка при делении двойного неравенства:
  при делении на отрицательное число оба знака неравенства меняют направление.
  В рассмотренном примере делитель $\pi>0$, поэтому знаки сохраняются.
\end{itemize}

\subsect{Короткий контроль перед сдачей}

\begin{enumerate}
  \item Аргументы приведены к удобному виду?
  \item После каждого существенного преобразования проверен общий множитель?
  \item При делении обосновано, что делитель не равен нулю?
  \item После группировки действительно появились одинаковые скобки?
  \item Период в формуле общего решения выбран верно?
  \item При отборе корней учтено условие $k\in\Z$?
  \item Все выбранные корни входят в заданный промежуток?
\end{enumerate}

\clearpage
\thispagestyle{fancy}

\sect{Тренировочный блок · Путь Б}

На каждый день предусмотрено по два задания.
Решения оформляйте полностью.
Если требуется отбор корней, сначала найдите общее решение уравнения,
затем выполните отбор.

\subsect{25.09}

\begin{enumerate}[label=\textbf{\arabic*.}]
  \item
  Решите уравнение
  \[
  \sin 2x-\sqrt3\,\sin x=0.
  \]

  \item
  Решите уравнение
  \[
  \sin x-\cos x=0.
  \]
\end{enumerate}

\subsect{26.09}

\begin{enumerate}[label=\textbf{\arabic*.},resume]
  \item
  Решите уравнение
  \[
  2\sin x+3\cos x=0.
  \]

  \item
  Решите уравнение
  \[
  \sin 2x-2\sin x+\cos x-1=0.
  \]
\end{enumerate}

\subsect{27.09}

\begin{enumerate}[label=\textbf{\arabic*.},resume]
  \item
  Решите уравнение
  \[
  \sin 2x+\sin x=0,
  \]
  затем отберите корни на промежутке
  \[
  \left[-\frac{\pi}{2};\frac{3\pi}{2}\right].
  \]

  \item
  Решите уравнение
  \[
  2\sin^2x+2\sin x\cos x+\sin x+\cos x=0,
  \]
  затем отберите корни на промежутке
  \[
  [0;2\pi).
  \]
\end{enumerate}

\clearpage
\thispagestyle{fancy}

\sect{Ответы}

\begin{table}[ht]
\centering
\begin{tabularx}{\linewidth}{@{}>{\bfseries}p{12mm}X@{}}
\toprule
№ & Ответ \\
\midrule

1 &
$x=\pi k$ или
$x=\pm\dfrac{\pi}{6}+2\pi k$,
$k\in\Z$.
\\

\addlinespace[2mm]

2 &
$x=\dfrac{\pi}{4}+\pi k$,
$k\in\Z$.
\\

\addlinespace[2mm]

3 &
$x=-\arctg\dfrac32+\pi k$,
$k\in\Z$.
\\

\addlinespace[2mm]

4 &
$x=2\pi k$ или
$x=-\dfrac{\pi}{6}+2\pi k$ или
$x=\dfrac{7\pi}{6}+2\pi k$,
$k\in\Z$.
\\

\addlinespace[2mm]

5 &
На заданном промежутке:
\[
x\in
\left\{
0,\,
\dfrac{2\pi}{3},\,
\pi,\,
\dfrac{4\pi}{3}
\right\}.
\]
\\

\addlinespace[2mm]

6 &
На заданном промежутке:
\[
x\in
\left\{
\dfrac{3\pi}{4},\,
\dfrac{7\pi}{6},\,
\dfrac{7\pi}{4},\,
\dfrac{11\pi}{6}
\right\}.
\]
\\

\bottomrule
\end{tabularx}
\end{table}

\vfill

{\small\color{softgray}
При проверке решений обращайте внимание на сохранность всех ветвей
после вынесения множителей, корректность периода и обоснованность
деления на тригонометрическое выражение.}

\end{document}